% GENERATED FILE. Edit canonical lessons, then run npm run generate:books.
% canonical-source-hash: fnv1a64:8b54169532648889
% canonical-lessons: JA-C390-kousokudouro, JA-C390-shasen, JA-C390-sokudo, JA-C390-menkyo, JA-C390-oikosu

\chapter{On the Expressway}
\label{ch:ja-on-the-expressway}
\hlchaptermodality{390}

\begin{chapteropening}
\textbf{By the end of this chapter:} \emph{I can talk about driving on the expressway.}

Complete the last lesson of chapter 390: I can talk about driving on the expressway.
\end{chapteropening}

\section[\texorpdfstring{\ja{こうそくどうろ} (kōsokudōro)}{こうそくどうろ (kōsokudōro)}]{\ja{こうそくどうろ} (kōsokudōro) — an expressway}

\label{lesson:JA-C390-kousokudouro}



\begin{quote}
Before the new one: say the Japanese for a footbridge (over a road), then the Japanese for to get through.
\end{quote}

\subsection*{You'll want to know: \ja{こうそくどうろ}}
\textbf{\ja{こうそくどうろ}} — \emph{kōsokudōro} — \textquotedblleft{}an expressway\textquotedblright{}.

\begin{quote}
\textbf{\ja{こうそくどうろで} \ja{まちに} \ja{いきました。}} — \emph{kōsokudōro de machi ni ikimashita} — \textquotedblleft{}We went into town on the expressway.\textquotedblright{}
\end{quote}

\begin{cousinweb}
\textbf{\ja{こうそくどうろ}} (\emph{kōsokudōro}) is Sino-Japanese (kango): four Chinese characters read with their on'yomi, readings borrowed from Chinese. They are \emph{kō}, \textquotedblleft{}high\textquotedblright{} (Middle Chinese \emph{kaw}), \emph{soku}, \textquotedblleft{}fast\textquotedblright{} (\emph{suwk}), \emph{dō}, \textquotedblleft{}street\textquotedblright{} (\emph{dawX}) and \emph{ro}, \textquotedblleft{}road\textquotedblright{} (\emph{luH}) (KANJIDIC2; Middle Chinese in Baxter's transcription).
\end{cousinweb}

\begin{grammarlens}[title={What you've built}]
One more word for getting around.
\end{grammarlens}

\subsection*{Your turn}
\begin{itemize}
\raggedright
  \item \emph{Say it:} \emph{kōsokudōro}
  \item \emph{Say it:} \emph{kōsokudōro}, once more
  \item \emph{Say it:} \emph{kōsokudōro de machi ni ikimashita} — We went into town on the expressway.
  \item \emph{Recall:} say the Japanese for a footbridge (over a road), then the Japanese for to get through, then say \emph{kōsokudōro} again
\end{itemize}

\begin{tcolorbox}[breakable,colback=teal!4,colframe=teal!35!black,title={Before you move on}]
What does \ja{こうそくどうろ} mean? (\textquotedblleft{}An expressway\textquotedblright{}.) Say it once more.
\end{tcolorbox}
\section[\texorpdfstring{\ja{しゃせん} (shasen)}{しゃせん (shasen)}]{\ja{しゃせん} (shasen) — a traffic lane}

\label{lesson:JA-C390-shasen}



\begin{quote}
Before the new one: say the Japanese for to get through, then the Japanese for an expressway.
\end{quote}

\subsection*{You'll want to know: \ja{しゃせん}}
\textbf{\ja{しゃせん}} — \emph{shasen} — \textquotedblleft{}a traffic lane\textquotedblright{}.

\begin{quote}
\textbf{\ja{この} \ja{どうろは} \ja{しゃせんが} \ja{みっつ} \ja{あります。}} — \emph{kono dōro wa shasen ga mittsu arimasu} — \textquotedblleft{}This road has three lanes.\textquotedblright{}
\end{quote}

\begin{cousinweb}
\textbf{\ja{しゃせん}} (\emph{shasen}) is Sino-Japanese (kango): two Chinese characters read with their on'yomi, readings borrowed from Chinese. They are \emph{sha}, \textquotedblleft{}car\textquotedblright{} (Middle Chinese \emph{tsyhae}) and \emph{sen}, \textquotedblleft{}line\textquotedblright{} (\emph{sjenH}) (KANJIDIC2; Middle Chinese in Baxter's transcription).
\end{cousinweb}

\begin{grammarlens}[title={What you've built}]
One more word for getting around.
\end{grammarlens}

\subsection*{Your turn}
\begin{itemize}
\raggedright
  \item \emph{Say it:} \emph{shasen}
  \item \emph{Say it:} \emph{shasen}, once more
  \item \emph{Say it:} \emph{kono dōro wa shasen ga mittsu arimasu} — This road has three lanes.
  \item \emph{Recall:} say the Japanese for to get through, then the Japanese for an expressway, then say \emph{shasen} again
\end{itemize}

\begin{tcolorbox}[breakable,colback=teal!4,colframe=teal!35!black,title={Before you move on}]
What does \ja{しゃせん} mean? (\textquotedblleft{}A traffic lane\textquotedblright{}.) Say it once more.
\end{tcolorbox}
\section[\texorpdfstring{\ja{そくど} (sokudo)}{そくど (sokudo)}]{\ja{そくど} (sokudo) — speed}

\label{lesson:JA-C390-sokudo}



\begin{quote}
Before the new one: say the Japanese for an expressway, then the Japanese for a traffic lane.
\end{quote}

\subsection*{You'll want to know: \ja{そくど}}
\textbf{\ja{そくど}} — \emph{sokudo} — \textquotedblleft{}speed\textquotedblright{}.

\begin{quote}
\textbf{\ja{あめの} \ja{ときは、そくどを} \ja{おとして} \ja{ください。}} — \emph{ame no toki wa, sokudo o otoshite kudasai} — \textquotedblleft{}Please slow down when it rains.\textquotedblright{}
\end{quote}

\begin{cousinweb}
\textbf{\ja{そくど}} (\emph{sokudo}) is Sino-Japanese (kango): two Chinese characters read with their on'yomi, readings borrowed from Chinese. They are \emph{soku}, \textquotedblleft{}fast\textquotedblright{} (Middle Chinese \emph{suwk}) and \emph{do}, \textquotedblleft{}degrees\textquotedblright{} (\emph{duH}) (KANJIDIC2; Middle Chinese in Baxter's transcription).
\end{cousinweb}

\begin{grammarlens}[title={What you've built}]
One more word for getting around.
\end{grammarlens}

\subsection*{Your turn}
\begin{itemize}
\raggedright
  \item \emph{Say it:} \emph{sokudo}
  \item \emph{Say it:} \emph{sokudo}, once more
  \item \emph{Say it:} \emph{ame no toki wa, sokudo o otoshite kudasai} — Please slow down when it rains.
  \item \emph{Recall:} say the Japanese for an expressway, then the Japanese for a traffic lane, then say \emph{sokudo} again
\end{itemize}

\begin{tcolorbox}[breakable,colback=teal!4,colframe=teal!35!black,title={Before you move on}]
What does \ja{そくど} mean? (\textquotedblleft{}Speed\textquotedblright{}.) Say it once more.
\end{tcolorbox}
\section[\texorpdfstring{\ja{めんきょ} (menkyo)}{めんきょ (menkyo)}]{\ja{めんきょ} (menkyo) — a licence; a driving licence}

\label{lesson:JA-C390-menkyo}



\begin{quote}
Before the new one: say the Japanese for a traffic lane, then the Japanese for speed.
\end{quote}

\subsection*{You'll want to know: \ja{めんきょ}}
\textbf{\ja{めんきょ}} — \emph{menkyo} — \textquotedblleft{}a licence; a driving licence\textquotedblright{}.

\begin{quote}
\textbf{\ja{あには} \ja{くるまの} \ja{めんきょを} \ja{もっています。}} — \emph{ani wa kuruma no menkyo o motte imasu} — \textquotedblleft{}My brother has a driving licence.\textquotedblright{}
\end{quote}

\begin{cousinweb}
\textbf{\ja{めんきょ}} (\emph{menkyo}) is Sino-Japanese (kango): two Chinese characters read with their on'yomi, readings borrowed from Chinese. They are \emph{men}, \textquotedblleft{}excuse\textquotedblright{} (Middle Chinese \emph{mjenX}) and \emph{kyo}, \textquotedblleft{}permit\textquotedblright{} (\emph{xjoX}) (KANJIDIC2; Middle Chinese in Baxter's transcription).
\end{cousinweb}

\begin{grammarlens}[title={What you've built}]
One more word for getting around.
\end{grammarlens}

\subsection*{Your turn}
\begin{itemize}
\raggedright
  \item \emph{Say it:} \emph{menkyo}
  \item \emph{Say it:} \emph{menkyo}, once more
  \item \emph{Say it:} \emph{ani wa kuruma no menkyo o motte imasu} — My brother has a driving licence.
  \item \emph{Recall:} say the Japanese for a traffic lane, then the Japanese for speed, then say \emph{menkyo} again
\end{itemize}

\begin{tcolorbox}[breakable,colback=teal!4,colframe=teal!35!black,title={Before you move on}]
What does \ja{めんきょ} mean? (\textquotedblleft{}A licence; a driving licence\textquotedblright{}.) Say it once more.
\end{tcolorbox}
\section[\texorpdfstring{\ja{おいこす} (oikosu)}{おいこす (oikosu)}]{\ja{おいこす} (oikosu) — to overtake; to pass}

\label{lesson:JA-C390-oikosu}



\begin{quote}
Before the new one: say the Japanese for speed, then the Japanese for a licence.
\end{quote}

\subsection*{You'll want to know: \ja{おいこす}}
\textbf{\ja{おいこす}} — \emph{oikosu} — \textquotedblleft{}to overtake; to pass\textquotedblright{}.

\begin{quote}
\textbf{\ja{こうそくどうろで} \ja{くるまを} \ja{おいこしました。}} — \emph{kōsokudōro de kuruma o oikoshimashita} — \textquotedblleft{}I overtook a car on the expressway.\textquotedblright{}
\end{quote}

Polite \textbf{\ja{おいこします}}, \emph{oikoshimasu}; a godan verb. Past (polite): \textbf{\ja{くるまが} \ja{じてんしゃを} \ja{おいこしました。}}, \emph{kuruma ga jitensha o oikoshimashita}.

\begin{cousinweb}
\textbf{\ja{おいこす}} (\emph{oikosu}) is a native Japanese word (wago), inherited rather than borrowed from Chinese, built from \emph{oi}, \textquotedblleft{}chasing\textquotedblright{} (from \emph{ou}) and \emph{kosu}, \textquotedblleft{}to cross over\textquotedblright{}; its kanji were attached later as kun'yomi (JMdict).
\end{cousinweb}

\begin{grammarlens}[title={What you've built}]
One more word for getting around.
\end{grammarlens}

\subsection*{Your turn}
\begin{itemize}
\raggedright
  \item \emph{Say it:} \emph{oikosu}
  \item \emph{Say it:} \emph{oikosu}, once more
  \item \emph{Say it:} \emph{kōsokudōro de kuruma o oikoshimashita} — I overtook a car on the expressway.
  \item \emph{Recall:} say the Japanese for speed, then the Japanese for a licence, then say \emph{oikosu} again
\end{itemize}

\begin{tcolorbox}[breakable,colback=teal!4,colframe=teal!35!black,title={Before you move on}]
What does \ja{おいこす} mean? (\textquotedblleft{}To overtake; to pass\textquotedblright{}.) Say it once more.
\end{tcolorbox}
